\section{Charts \& Tables}

\subsection*{[17.1] Terrain Effects Chart}

\begin{tcolorbox}[colback=cvsand!35,colframe=cvblue,boxrule=.6pt,arc=1pt]
\small
\begin{tabularx}{\textwidth}{@{}XYYY@{}}
  \toprule
  \textbf{Action Type} & \multicolumn{3}{c}{\textbf{Unit Type MP Cost}} \\
  \cmidrule(l){2-4}
   & \textbf{Foot} & \textbf{Light Horse} & \textbf{Heavy Horse} \\
  \midrule
  Enter Clear Hex & 1 & 1 & 1 \\
  Enter Hilltop Hex & 1 & 1 & 1 \\
  Enter Slope Hex & 1 & 1 & 1 \\
  Enter Woods Hex & 2 & P & P \\
  Enter Town Hex & P & P & P \\
  Enter Marsh Hex & 2 & 2 & 3 \\
  Enter Train Hex & 2 & 2 & 2 \\
  Enter Close Hex & 2 & 3 & 3 \\
  Enter Any Hex (except Town) via Lane Hexside & 1 & 1 & 1 \\
  Cross Stream Hexside & +2 & +2 & +3 \\
  Cross Ditch Hexside & +2 & +3 & +4 \\
  Cross Slope Hexside & \dash & +1 & +1 \\
  \bottomrule
\end{tabularx}

\vspace{1em}

\begin{tabularx}{\textwidth}{@{}XYYY@{}}
  \toprule
  \textbf{Action Type} & \multicolumn{3}{c}{\textbf{Attacking Effects}} \\
  \cmidrule(l){2-4}
   & \textbf{Foot} & \textbf{Light Horse} & \textbf{Heavy Horse} \\
  \midrule
  Attack into/out of Clear Hex & \dash & \dash & C \\
  Attack into/out of Hilltop Hex & \dash & \dash & C \\
  Attack into/out of Slope Hex & \dash & \dash & C \\
  Attack into/out of Train Hex & \dash & \dash & \dash \\
  Attack into/out of Close Hex & \dash & \dash & \dash \\
  Attack into Woods Hex & \dash & P & P \\
  Attack into Town Hex & P & P & P \\
  Attack out of Marsh Hex & H & H & H \\
  Attack Across Ditch Hexside & H & H & HC \\
  Attack Across Stream Hexside & H & H & HC \\
  Attack Up Across Slope Hexside & H & H & HC \\
  Attack Down Across Slope Hexside & \dash & \dash & \dash \\
  \bottomrule
\end{tabularx}

\smallskip
\textbf{Key:} \textbf{\#} = the MP cost to enter the hex. \textbf{+\#} = the additional cost to enter the hex through this type of hexside. \textbf{P} = entering or attacking into this hex type is prohibited to this unit type. \textbf{H} = this unit type's combat strength is halved when it is executing this action. \textbf{C} = Heavy Horse units can charge into, out of, or across this terrain type.
\end{tcolorbox}

\begin{tcolorbox}[colback=cvsand!55,colframe=cvred,boxrule=.6pt,arc=1pt,title={\sffamily\bfseries Course of Play},coltitle=cvblue]
\begin{multicols}{2}
\small
\begin{enumerate}[label=\textbf{Step \arabic*:},leftmargin=*,itemsep=.15em,topsep=0pt]
  \item Visibility Phase
  \item Allied Rally Phase
  \item Allied Artillery Phase
  \item Allied March Phase
  \item Allied Combat Phase
  \item Royalist Rally Phase
  \item Royalist Artillery Phase
  \item Royalist March Phase
  \item Royalist Combat Phase
  \item Game Turn Record Phase
\end{enumerate}
\end{multicols}
\end{tcolorbox}

\clearpage
\subsection*{[17.2] Artillery Table}

\begin{tcolorbox}[colback=cvsand!35,colframe=cvblue,boxrule=.6pt,arc=1pt]
\small
\begin{tabularx}{\textwidth}{@{}c*{6}{Y}@{}}
  \toprule
  \textbf{Die Roll} & \multicolumn{6}{c}{\textbf{Range in Hexes from Artillery Unit to Target}} \\
  \cmidrule(l){2-7}
    & \textbf{1} & \textbf{2} & \textbf{3} & \textbf{4} & \textbf{5} & \textbf{6+} \\
  \midrule
  1 & Dd & Dd & Dd & Dd & Dd & Dd \\
  2 & Dd & Dd & Dd & Dd & Dd & \dash \\
  3 & Dd & Dd & Dd & \dash & \dash & \dash \\
  4 & Dd & \dash & \dash & \dash & \dash & \dash \\
  5 & \dash & \dash & \dash & \dash & \dash & \dash \\
  6 & \dash & \dash & \dash & \dash & \dash & \dash \\
  \bottomrule
\end{tabularx}

\smallskip
\textbf{Key: Dd = Defender Disrupted:} The target unit is disrupted. \textbf{\dash{} = No Effect.} Range is always counted from the artillery unit's hex (exclusive) to the target hex (inclusive). Ranges less than 1 are not possible. Ranges greater than 6 are treated as 6. Disrupted units that suffer a Dd result due to artillery fire treat the result as if it were No Effect. A unit can never be eliminated as a direct result of artillery fire (though a unit disrupted by artillery fire can be eliminated due to a Dd result during combat).
\end{tcolorbox}

\subsection*{[17.3] Combat Table}

\begin{tcolorbox}[colback=cvsand!35,colframe=cvblue,boxrule=.6pt,arc=1pt]
\centering
\scriptsize
\resizebox{\textwidth}{!}{%
\begin{tabular}{@{}l*{11}{c}@{}}
  \toprule
  \textbf{Terrain} & \multicolumn{11}{c}{\textbf{Combat Ratio}} \\
  \midrule
  Marsh & 1--6 & 1--5 & 1--4 & 1--3 & 1--2 & 1--1 & 2--1 & 3--1 & 4--1 & 5--1 & 6--1 \\
  Clear/Slope/Hilltop & 1--5 & 1--4 & 1--3 & 1--2 & 1--1 & 2--1 & 3--1 & 4--1 & 5--1 & 6--1 & \\
  Train/Close---Horse & 1--4 & 1--3 & 1--2 & 1--1 & 2--1 & 3--1 & 4--1 & 5--1 & 6--1 & & \\
  Train/Woods---Foot & 1--3 & 1--2 & 1--1 & 2--1 & 3--1 & 4--1 & 5--1 & 6--1 & & & \\
  Close---Foot & 1--2 & 1--1 & 2--1 & 3--1 & 4--1 & 5--1 & 6--1 & & & & \\
  \midrule
  \textbf{Die Roll} & \multicolumn{11}{c}{\textbf{Result}} \\
  \midrule
  1 & Ad & \dash & \dash & Dx & Dd & Dd & Dd & De & De & De & De \\
  2 & Ad & Ad & \dash & \dash & Dx & Dd & Dd & Dd & De & De & De \\
  3 & Ae & Ad & Ad & \dash & \dash & Dx & Dd & Dd & Dd & De & De \\
  4 & Ae & Ad & Ad & Dx & \dash & \dash & Dx & Dd & Dd & Dd & De \\
  5 & Ae & Ae & Ad & Ad & Dx & \dash & \dash & Dx & Dd & Dd & Dd \\
  6 & Ae & Ae & Ae & Ad & Ad & Dx & \dash & \dash & Dx & Dd & Dd \\
  \bottomrule
\end{tabular}}

\raggedright\smallskip
\textbf{Key: Ad = Attacker Disrupted:} All attacking units are disrupted.
\textbf{Dd = Defender Disrupted:} All disrupted defending units are eliminated and all ordered defending units are disrupted.
\textbf{Dx = Disruption Exchange:} All disrupted defending units are eliminated; all ordered defending units are disrupted; attacking units whose printed combat strength is at least equal to the printed combat strengths of defending units are disrupted (attacker's choice of which units to disrupt).
\textbf{De = Defender Eliminated:} All defending units are immediately and permanently removed from play.
\textbf{Ae = Attacker Eliminated:} All attacking units are immediately and permanently removed from play.
\textbf{\dash{} = No Effect.} Attacks executed at a ratio less than the lowest ratio shown on a line are resolved at the lowest column shown on that line. Attacks executed at a ratio greater than 6--1 are resolved on the 6--1 column of that line.
\end{tcolorbox}
